\documentclass[a4paper,10pt]{article}
\usepackage[top=25mm,bottom=24mm,left=27mm,right=27mm,headheight=15pt,headsep=8mm,footskip=12mm]{geometry}
\usepackage{fontspec,amsmath,amsthm,unicode-math}
\usepackage[spanish,es-nodecimaldot,es-noquoting]{babel}
\usepackage{microtype,xcolor,fancyhdr,titlesec,booktabs,longtable,array,tabularx}
\usepackage{graphicx,tikz,enumitem,float,needspace}
\usetikzlibrary{arrows.meta,calc,positioning,decorations.pathreplacing,matrix}
\usepackage[unicode,hidelinks]{hyperref}
\usepackage[nameinlink,noabbrev]{cleveref}
\setmainfont{STIXTwoText-Regular.otf}[BoldFont=STIXTwoText-Bold.otf,ItalicFont=STIXTwoText-Italic.otf,BoldItalicFont=STIXTwoText-BoldItalic.otf]
\setmathfont{STIXTwoMath-Regular.otf}
\hypersetup{pdftitle={Moonshine, dualidad T y teoría M desde la estructura discreta del continuo},pdfauthor={Oumar Haidara Fall; Rubén Ramos Balsa},pdfsubject={Artículo IV. Borrador de trabajo para revisión académica}}
\newcommand{\Autores}{Oumar Haidara Fall \textperiodcentered\ Rubén Ramos Balsa}
\newcommand{\Titulo}{Moonshine, dualidad T y teoría M desde la estructura discreta del continuo}
\newcommand{\RR}{\mathbb R}\newcommand{\ZZ}{\mathbb Z}\newcommand{\QQ}{\mathbb Q}\newcommand{\CC}{\mathbb C}\newcommand{\FF}{\mathbb F}
\newcommand{\Id}{\operatorname{id}}\newcommand{\tr}{\operatorname{tr}}
\newcommand{\Span}{\operatorname{span}}\newcommand{\Cyl}{\operatorname{Cyl}}
\newcommand{\square}{\mdlgwhtsquare}
\newcommand{\TPK}{\mathrm{TPK}}\newcommand{\APP}{\mathrm{APP}}\newcommand{\TRIT}{\{-1,0,+1\}}
\newtheorem{theorem}{Teorema}[section]
\newtheorem{lemma}[theorem]{Lema}
\newtheorem{proposition}[theorem]{Proposición}
\newtheorem{corollary}[theorem]{Corolario}
\theoremstyle{definition}
\newtheorem{definition}[theorem]{Definición}
\newtheorem{axiom}[theorem]{Axioma}
\theoremstyle{remark}\newtheorem{remark}[theorem]{Observación}
\numberwithin{equation}{section}
\setlength{\parindent}{1em}\setlength{\parskip}{3pt plus 1pt minus .5pt}
\setlength{\emergencystretch}{2em}
\renewcommand{\normalsize}{\fontsize{10.5}{13.4}\selectfont}\normalsize
\setlist{nosep,topsep=5pt}
\titleformat{\section}{\fontsize{13}{16}\selectfont\bfseries}{\thesection}{.65em}{}
\titleformat{\subsection}{\fontsize{11.2}{14}\selectfont\bfseries}{\thesubsection}{.65em}{}
\titleformat{\subsubsection}{\normalsize\bfseries}{\thesubsubsection}{.65em}{}
\AddToHook{cmd/subsection/before}{\Needspace{7\baselineskip}}
\AddToHook{cmd/subsubsection/before}{\Needspace{6\baselineskip}}
\setcounter{tocdepth}{2}
\titlespacing*{\section}{0pt}{17pt}{9pt}\titlespacing*{\subsection}{0pt}{13pt}{6pt}
\pagestyle{fancy}\fancyhf{}
\fancyhead[L]{\fontsize{8}{10}\selectfont Moonshine, dualidad T y pantallas dimensionales}
\fancyhead[R]{\fontsize{8}{10}\selectfont Artículo IV}
\fancyfoot[C]{\thepage}\renewcommand{\headrulewidth}{.2pt}
\raggedbottom
\begin{document}
\begin{center}
{\fontsize{18}{22}\selectfont\bfseries\Titulo\par}
\vspace{10mm}
{\large\Autores\par}
\vspace{5mm}
{\small Artículo IV de la serie HMT--MD\par}
\vspace{2mm}
{\small Borrador de trabajo para revisión académica\par}
{\footnotesize Integración y recopilación documental generadas por ChatGPT Astra\par}
\end{center}
\vspace{6mm}
\begin{center}\textbf{Resumen}\end{center}
Se desarrolla la prolongación excepcional y dimensional de un núcleo formal
matemático denominado \emph{Holografía Modular Triádica}, en adelante HMT.
La construcción comienza en el soporte aritmético \emph{All Possible Paths}
(APP), incorpora la unidad ternaria orientada TRIT y compone sus operaciones
mediante el \emph{Topological Phase Kernel} (TPK). El estado resultante
conserva conjuntamente orientación, hojas aritméticas, residuos, cocientes,
frontera y memoria. Su prolongación compatible suministra la fuente común
de la realización numérica y de la incidencia excepcional.

El registro dodecafásico \(K\), construido dentro de esa cadena, cumple dos
funciones que se desarrollan mediante mapas explícitos. Su incidencia
interviene en la selección de una bandera de Witt y en la construcción
reticular que alcanza Leech y el álgebra de operadores de vértice
\(V^\natural\). Su proyección sobre una componente irreducible de la
representación de \(A_5\) determina una dirección no nula y la reducción
ortogonal \(12\to11\to10\). La construcción excepcional se acompaña de
los automorfismos de los sectores y de las realizaciones de orbifold de
órdenes \(2\) y \(3\), manteniendo identificados sus dominios.

La fase nonádica admite una representación de Weyl a cada profundidad.
Las inclusiones que conservan la fibra completa producen un álgebra
uniformemente hiperfinita y su realización tracial de tipo \(II_1\).
El intercambio de dos coordenadas enteras y la inversión del radio
determinan una dualidad exacta, tanto para la forma cuadrática compacta
como para su Hamiltoniano autoadjunto. Se prueba asimismo su formulación
covariante bajo el cambio de sección residual. Las pantallas reales,
combinatorias y reticulares se coordinan por sus morfismos, incluida la
reducción isotrópica \(26\to24\).

Finalmente, la composición de trayectorias con frontera conservada
define un cociclo de acción. La realización undecadimensional se formula
con sus campos, dominios operatorios y entrelazamientos: las identidades
algebraicas demostradas y el teorema de transporte de la supercarga
determinan con precisión el alcance de la interfaz física. El artículo
reúne en su propio cuerpo el núcleo común, la generación de los registros
y el corredor excepcional que necesitan estas construcciones.

\medskip
\noindent\textbf{Palabras clave:} Holografía Modular Triádica; estructura
discreta del continuo; registro dodecafásico; simetrías excepcionales;
Moonshine; dualidad T; representaciones de Weyl; pantallas dimensionales;
teoría M.

\clearpage
\section{Introducción}
\label{iv:introduccion}

La relación entre una estructura matemática y su realización física
requiere conocer qué operaciones producen el objeto, qué información
conservan sus representaciones y qué ecuaciones se transportan entre
ellas. Una dimensión, un valor espectral o un grupo de automorfismos
constituyen resultados distintos de esa investigación. Cuando proceden
de un mismo estado, su relación se expresa mediante los mapas que los
producen. Este artículo estudia esa relación en la prolongación de HMT
hacia la incidencia excepcional, las realizaciones hilbertianas y los
sistemas de compactificación.

La tesis que organiza la serie sitúa la estructura discreta del continuo
antes de su evaluación real. APP proporciona las dos evaluaciones
aritméticas sobre un soporte común; TRIT conserva los regímenes
orientados; TPK selecciona, transporta y prolonga los estados con sus
datos de frontera y memoria. La reducción de cilindros compatibles
determina coordenadas numéricas. La proyección excepcional conserva la
incidencia que se realiza mediante códigos, retículos y álgebras
graduadas. El valor y la estructura que lo genera permanecen enlazados
por esta procedencia, aun cuando cada representación retenga información
diferente.

La función del registro \(K\) permite seguir ese enlace de forma
especialmente concreta. Primero se construyen sus doce componentes
desde el registro firmado y las operaciones del TPK. Su lectura
racional conserva el orden de los bloques con base y origen de fase
fijados. La incidencia de esos componentes interviene después en la
bandera que conduce del diseño de Witt al retículo de Leech.
En otra representación declarada, el mismo registro produce la
dirección \(u_K\) que distingue una recta dentro del espacio
\(V_{11}\). La reducción a \(V_{10}\) tiene entonces un núcleo
especificado, una norma exacta y una ley de covariancia. Este recorrido
explica por qué es indispensable conservar la construcción de \(K\)
antes de utilizar su valor o su dirección.

El corredor excepcional se expone hasta \(V^\natural\), con la
construcción de Frenkel--Lepowsky--Meurman y los resultados de Moonshine
en sus dominios propios~\cite{flm,borcherds}. La elección del retículo,
la composición de sus sectores y la acción de sus automorfismos ocupan
lugares distintos dentro de la prueba. Las realizaciones de órdenes
\(2\) y \(3\) se incorporan con sus referencias y transformaciones.
La función modular se relaciona posteriormente con la inversión del
radio mediante una identidad explícita en el eje imaginario del
semiplano superior. Esa identidad permite componer una relación precisa
entre parámetros sin sustituir los operadores de cada representación.

La realización hilbertiana aporta otro paso necesario. En cada
profundidad se conserva la etiqueta de fase de una base ortonormal.
El desplazamiento y el carácter de fase generan toda el álgebra
matricial de esa fibra. Al refinar, la inclusión mediante la identidad
de la fibra nueva conserva la unidad y la traza normalizada; fijar una
sola continuación produce una inclusión en una esquina. Las pruebas
del cierre inductivo y de la realización tracial hacen visible la
consecuencia de esa elección. La conservación de la unidad deja así
de ser sólo una formulación conceptual: interviene en la acción de
las inclusiones y en el tipo del álgebra resultante.

La dualidad compacta se desarrolla a continuación sobre un dominio
estable. El intercambio de las coordenadas enteras exige transportar
también el radio y, cuando cambia la sección residual, la relación de
equivalencia que conserva el entero. El artículo demuestra la
involución, la invariancia de la forma cuadrática, la autoadjunción
del Hamiltoniano y la igualdad de sus dominios bajo conjugación
unitaria. La evaluación sobre clases incluye el cociente euclídeo;
su mínimo puede calcularse mediante dos candidatos enteros. Así se
conserva el contenido aritmético del cambio de representante durante
la construcción de la equivalencia espectral.

Las pantallas dimensionales reúnen estos resultados con la incidencia.
La perforación de un código, la eliminación ortogonal de una dirección
y el cociente isotrópico de un retículo son operaciones diferentes.
Cada una determina un rango y conserva un núcleo propio. Su
coordinación utiliza esas presentaciones completas. La realización
undecadimensional recibe después las pantallas y la dualidad mediante
un mapa que transporta campos, graduación, acción y condiciones de
frontera. La relación de supercarga se expresa como igualdad de
operadores con sus dominios, y su teorema de transporte prueba la
conclusión sobre la imagen de la isometría declarada.

El orden expositivo mantiene estas dependencias. Se reproduce primero
el núcleo formal común y su desarrollo, después la generación
coinductiva, el registro \(K\), las dos construcciones de \(\alpha\)
y la incidencia excepcional. Sobre esa base se introducen la
realización hilbertiana, la dualidad, las pantallas y el cociclo de
acción. La repetición del núcleo entre artículos garantiza que este
texto pueda leerse por separado; las fuentes del tratado
integral~\cite{hmtedicionintegral} documentan la procedencia de las
construcciones reunidas. Las demostraciones necesarias se sitúan en
el propio desarrollo y las referencias internas permiten recorrerlas
en ambos sentidos.

\clearpage
\tableofcontents
\clearpage
\input{sections/nucleo.tex}
\clearpage
\input{sections/extension.tex}
\clearpage
\input{sections/generacion.tex}
\clearpage
\input{sections/registro_k.tex}
\clearpage
\input{sections/alpha.tex}
\clearpage
\input{sections/excepcional.tex}
\clearpage
\input{sections/01_hilbert_nonadico.tex}
\clearpage
% Procedencia: integral 2249, cap. 24, pp. 609--618; capitulo_24_dualidad_t.tex.
% Las fórmulas de dominio operatorio se explicitan sobre la realización ya construida.
\section{Dualidad T: coordenadas enteras, cambio de sección y equivalencia operatoria}
\label{iv:dualidad}

La descomposición aritmética de APP conserva el representante residual
y su cociente. La evolución TPK transporta además su orientación y
la memoria de los cruces de frontera. El intercambio de las dos
coordenadas exige conservar esta distinción y precisar el dominio
en que la operación vuelve a estar definida. Sobre ese dominio se
construye primero una involución aritmética; su forma cuadrática
admite después una representación por operadores autoadjuntos.

\subsection{Dominio estable del intercambio}

En la sección residual positiva,
\begin{equation}
 r_9^+(n)=1+((n-1)\bmod9),\qquad
 Q_9^+(n)=\frac{n-r_9^+(n)}9,
 \qquad n=r_9^+(n)+9Q_9^+(n).
 \label{iv:eq:division-positiva}
\end{equation}
Un intercambio puede situar un cociente arbitrario en la primera
coordenada. Por ello su dominio estable es
\(\mathscr D_0=\ZZ^2\times\RR_{>0}\), no solamente la sección
\(1\leq r\leq9\). La prolongación del dominio permite definir
una operación total; la admisibilidad de una pareja como lectura de
una ruta TPK conserva por separado sus reglas de selección.

\begin{definition}[Forma compacta normalizada e intercambio]
\label{iv:def:dualidad}
Para \(d=(m,w,R_0)\in\mathscr D_0\), sean
\begin{equation}
 E_{R_0}(m,w)=\frac{m^2}{R_0^2}+w^2R_0^2,
 \qquad
 \mathsf T_0(m,w,R_0)=(w,m,R_0^{-1}).
 \label{iv:eq:energia-t}
\end{equation}
Los símbolos \(m,w\) nombran las dos coordenadas enteras
transportadas; \(R_0\) es la escala positiva adimensional de
la forma. Su interpretación mediante momento, enrollamiento y
radio físico se establece por un mapa de realización posterior.
\end{definition}

\begin{theorem}[Involución y conservación de la forma cuadrática]
\label{iv:thm:dualidad-escalar}
La aplicación \(\mathsf T_0\) es una involución de
\(\mathscr D_0\), y
\[
 E_{R_0^{-1}}(w,m)=E_{R_0}(m,w).
\]
\end{theorem}
\begin{proof}
Dos intercambios restituyen \((m,w)\), y dos inversiones restituyen
\(R_0\). Además,
\[
 E_{R_0^{-1}}(w,m)=w^2R_0^2+\frac{m^2}{R_0^2}
 =E_{R_0}(m,w).
\]
\end{proof}

La equivalencia vale para cada radio y su recíproco. En \(R_0=1\),
el radio queda fijo y la operación actúa sobre una misma forma.
Una pareja concreta es fija por la involución completa precisamente
cuando, además, \(m=w\).

\subsection{Hamiltoniano compacto y dominio autoadjunto}

Sea \(\mathcal H_{\mathrm{lat}}=\ell^2(\ZZ^2)\), con base
ortonormal \(\{|m,w\rangle\}\). Definimos
\begin{align}
 H(R_0)|m,w\rangle&=E_{R_0}(m,w)|m,w\rangle,\label{iv:eq:hamiltoniano}\\
 D(H(R_0))&=\left\{\psi:\sum_{m,w\in\ZZ}
 E_{R_0}(m,w)^2|\psi_{m,w}|^2<\infty\right\}.\label{iv:eq:dominio-h}
\end{align}
Los vectores de soporte finito pertenecen al dominio y forman un
núcleo del operador diagonal. La forma cuadrática asociada tiene
dominio
\[
 D(H(R_0)^{1/2})=\left\{\psi:\sum_{m,w}
 E_{R_0}(m,w)|\psi_{m,w}|^2<\infty\right\}.
\]

\begin{proposition}[Autoadjunción y resolvente compacto]
\label{iv:prop:h-autoadjunto}
El operador \(H(R_0)\) es autoadjunto y no negativo. Sus
autovalores son \(E_{R_0}(m,w)\), contados con multiplicidad,
y \((H(R_0)+I)^{-1}\) es compacto.
\end{proposition}
\begin{proof}
Un operador de multiplicación por una sucesión real es simétrico
en \eqref{iv:eq:dominio-h}. Si \(\eta\) pertenece al dominio
del adjunto, probar la identidad del adjunto contra cada vector
\(|m,w\rangle\) obliga a que su imagen tenga coordenadas
\(E_{R_0}(m,w)\eta_{m,w}\). La condición de pertenecer a
\(\ell^2\) es exactamente \eqref{iv:eq:dominio-h}; por tanto
el adjunto tiene el mismo dominio y la misma acción. La no
negatividad se obtiene sumando los pesos no negativos.

Con \(c_{R_0}=\min(R_0^{-2},R_0^2)>0\), se tiene
\(E_{R_0}(m,w)\geq c_{R_0}(m^2+w^2)\). Fuera de cualquier
conjunto finito creciente, los coeficientes
\((1+E_{R_0}(m,w))^{-1}\) tienden uniformemente a cero.
Truncar el resolvente en dichos conjuntos produce operadores
de rango finito que convergen en norma, lo cual prueba su
compacidad. La base diagonal determina el espectro anunciado.
\end{proof}

\begin{theorem}[Conjugación unitaria, incluidos los dominios]
\label{iv:thm:t-unitaria}
El intercambio
\(U_T|m,w\rangle=|w,m\rangle\) es unitario,
\(U_T^2=I\), y satisface
\begin{equation}
 U_TD(H(R_0))=D(H(R_0^{-1})),
 \qquad U_TH(R_0)U_T^{-1}=H(R_0^{-1}).
 \label{iv:eq:conjugacion-h}
\end{equation}
También transporta los dominios de las formas cuadráticas.
\end{theorem}
\begin{proof}
El intercambio es una permutación involutiva de la base
ortonormal. Para una sucesión \(\psi\), la identidad del
teorema \ref{iv:thm:dualidad-escalar} da
\[
 \sum_{m,w}E_{R_0^{-1}}(m,w)^2|(U_T\psi)_{m,w}|^2
 =\sum_{m,w}E_{R_0}(m,w)^2|\psi_{m,w}|^2.
\]
Así se obtiene la igualdad de dominios. Sobre cada coordenada,
\(U_TH(R_0)U_T^{-1}\) multiplica por
\(E_{R_0}(w,m)=E_{R_0^{-1}}(m,w)\). La misma suma, con
la primera potencia del peso, demuestra el resultado para las
formas cuadráticas.
\end{proof}

La igualdad escalar y la igualdad operatoria son dos realizaciones
de la misma involución. La segunda conserva los dominios que
permiten interpretar la primera como equivalencia espectral,
no sólo como una igualdad sobre vectores de soporte finito.

\subsection{Cambio de representante y memoria del cociente}

La sección cero usa \(\bar r\in\{0,\ldots,8\}\). El cambio
que conserva el entero de \eqref{iv:eq:division-positiva} es
\begin{equation}
 C(r,Q)=(\bar r,Q+\mathbf1_{\{r=9\}}).
 \label{iv:eq:cambio-seccion}
\end{equation}
La modificación del cociente forma parte de la operación.
Para \(n=9\), las dos parejas son \((9,0)\) y \((0,1)\).
Sus energías no corregidas son \(81/R_0^2\) y \(R_0^2\).
Representar el mismo entero no implica, por tanto, igualdad de
esas dos inserciones en una forma que distingue las coordenadas.

La relación de equivalencia se conserva definiendo
\[
 L_9=\ZZ(9,-1),\qquad S(x_1,x_2)=(x_2,x_1),
 \qquad q_{R_0}(x)=\frac{x_1^2}{R_0^2}+x_2^2R_0^2.
\]
Para \(x\in\ZZ^2\), sea
\begin{equation}
 \mathcal E_{R_0}^{L_9}([x])
 =\min_{k\in\ZZ}q_{R_0}(x+k(9,-1)).
 \label{iv:eq:energia-cociente}
\end{equation}
Esta construcción opera sobre clases. La memoria que distingue
representantes individuales permanece en el estado TPK anterior
al cociente; la función de \eqref{iv:eq:energia-cociente}
determina la evaluación invariante de la clase elegida.

\begin{theorem}[Dualidad covariante de las secciones]
\label{iv:thm:dualidad-cociente}
La energía de \eqref{iv:eq:energia-cociente} está bien definida.
La transformación
\[
 \mathsf T_{\mathrm{cov}}([x]_{L_9},R_0)
 =([Sx]_{SL_9},R_0^{-1})
\]
es una biyección entre
\((\ZZ^2/L_9)\times\RR_{>0}\) y
\((\ZZ^2/SL_9)\times\RR_{>0}\), con inversa dada por el
segundo intercambio, y
\[
 \mathcal E_{R_0^{-1}}^{SL_9}([Sx])
 =\mathcal E_{R_0}^{L_9}([x]).
\]
\end{theorem}
\begin{proof}
Reemplazar \(x\) por \(x+k_0(9,-1)\) reindexa los candidatos
al mínimo. Éste existe: la función de \(k\) es cuadrática y
tiene coeficiente principal \(81/R_0^2+R_0^2>0\). El cambio
\((9,Q)\mapsto(0,Q+1)\) tiene diferencia en \(L_9\), por
lo que representa una misma clase. El intercambio transporta
la relación de equivalencia a \(SL_9\), y
\[
 \min_{\ell\in L_9}q_{R_0^{-1}}(Sx+S\ell)
 =\min_{\ell\in L_9}q_{R_0}(x+\ell).
\]
Finalmente, \(S^2=I\) recupera el cociente y el radio iniciales.
\end{proof}

La evaluación del mínimo puede efectuarse exactamente con dos
candidatos. Si \(x=(m,w)\), el mínimo real de la cuadrática
ocurre en
\[
 k_* =\frac{wR_0^4-9m}{81+R_0^4}.
\]
El mínimo entero se alcanza en \(\lfloor k_*\rfloor\) o en
\(\lceil k_*\rceil\). Esta fórmula no altera la equivalencia;
explicita su evaluación finita y permite comprobar el efecto
del cambio de sección sin descartar el cociente.

\subsection{Normalización de radio en la realización de cuerda}

Sea \(\alpha'>0\) una escala de la realización de cuerda.
Este símbolo designa la escala de cuerda y se mantiene distinto
de la constante de estructura fina \(\alpha\). Definimos
\begin{equation}
 \mathsf T_{\alpha'}(m,w,R)
 =\left(w,m,\frac{\alpha'}R\right),\qquad
 \mathcal N_{\alpha'}(m,w,R)
 =\left(m,w,\frac R{\sqrt{\alpha'}}\right).
 \label{iv:eq:normalizacion-cuerda}
\end{equation}

\begin{theorem}[Conjugación exacta por la escala]
\label{iv:thm:normalizacion-cuerda}
La normalización satisface
\[
 \mathcal N_{\alpha'}\mathsf T_{\alpha'}
 =\mathsf T_0\mathcal N_{\alpha'}.
\]
Si
\[
 E_{\mathrm{str}}(m,w,R)
 =\frac{m^2}{R^2}+\frac{w^2R^2}{\alpha'^2},
\]
entonces
\(E_{\mathrm{str}}=\alpha'^{-1}E\circ\mathcal N_{\alpha'}\).
Las combinaciones
\[
 p_L=\frac mR+\frac{wR}{\alpha'},\qquad
 p_R=\frac mR-\frac{wR}{\alpha'}
\]
se transforman mediante \(p_L\mapsto p_L\) y \(p_R\mapsto-p_R\).
\end{theorem}
\begin{proof}
Las dos composiciones de la primera igualdad dan
\((w,m,\sqrt{\alpha'}/R)\). Por otra parte,
\[
 E_{R/\sqrt{\alpha'}}(m,w)
 =\frac{\alpha'm^2}{R^2}+\frac{w^2R^2}{\alpha'}
 =\alpha'E_{\mathrm{str}}(m,w,R).
\]
Sustituir \((w,m,\alpha'/R)\) en las expresiones de
\(p_L,p_R\) da, respectivamente, las dos transformaciones
afirmadas.
\end{proof}

La conjugación fija qué significa expresar el mismo resultado en
radio normalizado o físico. La asignación de una tensión y una
unidad a \(\alpha'\) pertenece al mapa de realización y tiene
su propio origen; no es un dato que intervenga en la producción
de los registros APP--TRIT--TPK.

\subsection{Relación con la transformación modular}

La inversión del radio tiene una representación precisa en el
eje imaginario del semiplano superior. Sean
\(\iota(R_0)=\mathrm iR_0^2\) y
\(\mathsf S(\tau)=-1/\tau\). Entonces
\begin{equation}
 \mathsf S\circ\iota(R_0)=\iota(R_0^{-1}).
 \label{iv:eq:semiconjugacion-modular}
\end{equation}
La igualdad resulta de
\(-1/(\mathrm iR_0^2)=\mathrm iR_0^{-2}\). Para la función
modular \(J=j-744\), cuya invariancia se conserva en el
desarrollo de Moonshine, se deduce
\(J(\mathrm iR_0^2)=J(\mathrm iR_0^{-2})\).

El mapa \(\iota\) expresa una semiconjugación de parámetros.
El intercambio \(U_T\) actúa en las coordenadas de
\(\ell^2(\ZZ^2)\); \(\mathsf S\) actúa en el semiplano
superior; el automorfismo dual de un orbifold actúa en sus
sectores graduados. La escritura de sus dominios permite
componer relaciones específicas entre ellos conservando la
información de cada representación. Los resultados aquí
reunidos proceden del desarrollo de dualidad del integral,
ahora con dominios operatorios explícitos~\cite{hmtedicionintegral}.

\clearpage
% Procedencia: integral 2249, cap. 25, pp. 619--640, especialmente §25.13.
% Requiere la incorporación material de exc:k-direccion, exc:leech y el núcleo común.
\section{Pantallas dimensionales y reducción coordinada del estado}
\label{iv:pantallas}

Los rangos que aparecen en la prolongación excepcional y en la
compactificación pertenecen a objetos diferentes. Para relacionarlos
se conserva su presentación: generadores, relaciones, producto
interno y dirección seleccionada. La dimensión es entonces una
consecuencia del módulo de canales y de sus relaciones. Su
transporte se formula mediante mapas que actúan sobre ese módulo,
además de actuar sobre el estado que lo determina.

\subsection{Canales, relaciones y dimensión efectiva}

Sea \(\mathcal S_n\) el conjunto de estados supervivientes de
profundidad \(n\), con restricciones compatibles
\(r_{m,n}:\mathcal S_m\to\mathcal S_n\), \(m\geq n\).
El dato de frontera \(\beta_n(x)\) conserva la información
necesaria para decidir la compatibilidad de una prolongación:
hoja, orientación, cociente, memoria y clase de incidencia según
el lector considerado. Las restricciones satisfacen
\[
 r_{\ell,n}=r_{m,n}r_{\ell,m},\qquad
 \beta_n r_{m,n}=b_{m,n}\beta_m.
\]
Estas igualdades son las condiciones de naturalidad de la
presentación y mantienen los mismos datos al restringir por
etapas o directamente.

\begin{definition}[Pantalla de canales]
\label{iv:def:pantalla}
Fijado un cuerpo \(\mathbb K_n\), sea \(M_n(x)\) el espacio
generado por los canales de prolongación y \(\mathcal R_n(x)\)
el subespacio de relaciones. La pantalla es
\[
 \Sigma_n(x)=\bigl(\mathbb K_n,M_n(x),\beta_n(x),
                    \mathcal R_n(x)\bigr),
\]
y su dimensión efectiva es
\begin{equation}
 d_{\Sigma_n}(x)=\dim_{\mathbb K_n}
                    \bigl(M_n(x)/\mathcal R_n(x)\bigr).
 \label{iv:eq:dimension-pantalla}
\end{equation}
Si \(M_n(x)\cong\mathbb K_n^{g_n}\) y las columnas de
\(R_n(x)\) generan las relaciones, entonces
\(d_{\Sigma_n}(x)=g_n-\operatorname{rank}R_n(x)\).
\end{definition}

\begin{proposition}[Invariancia de presentación y transporte]
\label{iv:prop:pantalla-transporte}
La dimensión \eqref{iv:eq:dimension-pantalla} se conserva
bajo cambios de bases invertibles. Para dos pantallas sobre un mismo cuerpo
\(\mathbb K_n=\mathbb K_m=\mathbb K\), si la aplicación lineal
\(F_\gamma:M_n(x)\to M_m(T_\gamma x)\) envía
\(\mathcal R_n(x)\) en \(\mathcal R_m(T_\gamma x)\), induce
de manera única una aplicación lineal entre los cocientes de pantalla.
\end{proposition}
\begin{proof}
Un cambio de generadores y de relaciones transforma \(R\) en
\(URV\), con \(U,V\) invertibles, y conserva su rango. Para
el descenso, definimos
\(\overline F_\gamma([v])=[F_\gamma(v)]\). Si dos
representantes difieren en una relación, la linealidad implica que sus imágenes difieren
en una relación de la pantalla terminal; la aplicación está
bien definida. La proyección al cociente es sobreyectiva y
garantiza la unicidad.
\end{proof}

El cambio dimensional durante el transporte queda determinado
por
\[
 \Delta d(\gamma;x)=d_{\Sigma_m}(T_\gamma x)
                         -d_{\Sigma_n}(x).
\]
Esta expresión conserva el estado y su presentación de canales.
Un rango aislado no permite reconstruir ni las relaciones ni
el mapa que lo produjo.

\subsection{Perforación ternaria y pantalla de síndromes}

El código \(C_W=[12,6,6]_3\), construido a partir de la
incidencia excepcional, produce otra pantalla mediante
perforación de una coordenada. Denotemos por
\(p:C_W\to\FF_3^{11}\) ese mapa y por
\(G_{11}=p(C_W)\) su imagen.

\begin{proposition}[Código perforado y partición perfecta]
\label{iv:prop:codigo-once}
El mapa \(p\) es inyectivo. El código \(G_{11}\) tiene
parámetros \([11,6,5]_3\), y sus bolas de Hamming de radio
\(2\) particionan \(\FF_3^{11}\).
\end{proposition}
\begin{proof}
Un vector del núcleo de \(p\) tendría soporte contenido en
la única coordenada retirada. La distancia mínima \(6\)
de \(C_W\) excluye todo vector no nulo de esa forma.
Así, la dimensión sigue siendo \(6\). Perforar disminuye
el peso a lo sumo en una unidad, luego la distancia mínima
de la imagen es al menos \(5\). Las bolas de radio \(2\)
son por ello disjuntas. Cada una contiene
\[
 1+2\binom{11}{1}+2^2\binom{11}{2}
 =1+22+220=243
\]
vectores. Como \(|G_{11}|=3^6\) y
\(3^6\cdot243=3^{11}\), las bolas cubren el ambiente.
Para comprobar que la distancia es exactamente \(5\),
tómese un vector de peso \(3\). Su bola decodificadora
no puede tener centro cero y posee un centro de peso a lo
sumo \(3+2=5\); dicho centro debe tener peso al menos
\(5\). Se obtiene así una palabra de peso \(5\).
\end{proof}

El entero \(243=3^{11-6}\) es aquí el número de clases
del cociente de síndromes \(\FF_3^{11}/G_{11}\). Su
función se fija por la perforación y la distancia del
código; la igualdad de cardinal con otra fibra no sustituye
el mapa entre ambas. Esta pantalla combinatoria y la
siguiente pantalla real conservan sus operaciones propias.

\subsection{La dirección dodecafásica y la reducción de rango}

La construcción de \ref{exc:k-direccion} parte del registro
\(K\) ya obtenido por APP--TRIT--TPK y de la carta de
representación declarada sobre \(\RR^{12}\). En esa carta,
\begin{equation}
 V_{11}=\mathbf1^\perp,\qquad
 P_{11}=I_{12}-\frac1{12}J_{12},\qquad
 u_K=P_3P_{11}K.
 \label{iv:eq:p11-uk}
\end{equation}
Aquí \(J_{12}\) es la matriz de unos y \(P_3\) es el
proyector isótipo especificado por la acción de \(A_5\).
La suma finita que define \(P_3\), sus representantes y la
generación de \(K\) se conservan en la sección citada de
este mismo artículo. En particular,
\[
 \|u_K\|^2=\frac{6638585+2275584\sqrt5}{20}>0.
\]
La norma positiva define \(\ell_K=\RR u_K\), y el
proyector
\begin{equation}
 P_{10}=P_{11}-\frac{u_Ku_K^{\mathsf T}}{\|u_K\|^2}
 \label{iv:eq:p10}
\end{equation}
tiene imagen \(V_{10}=V_{11}\cap u_K^\perp\).

\begin{proposition}[Dos reducciones ortogonales]
\label{iv:prop:12-11-10}
Se cumplen
\[
 \RR^{12}=\RR\mathbf1\mathbin{\oplus^\perp}V_{11},\qquad
 V_{11}=\ell_K\mathbin{\oplus^\perp}V_{10},\qquad
 \dim V_{10}=10.
\]
Además, \(P_{10}P_{11}=P_{10}\), \(P_{10}u_K=0\), y
\(P_{10}\) es un proyector ortogonal.
\end{proposition}
\begin{proof}
El operador \(J_{12}/12\) proyecta sobre \(\RR\mathbf1\).
Sea \(E_K=u_Ku_K^{\mathsf T}/\|u_K\|^2\). Es simétrico,
de rango uno, y \(E_K^2=E_K\). Como \(u_K\in V_{11}\),
\(P_{11}E_K=E_KP_{11}=E_K\). Por tanto,
\[
 (P_{11}-E_K)^2=P_{11}-2E_K+E_K=P_{11}-E_K.
\]
Su imagen es el complemento ortogonal de \(\ell_K\) dentro
de \(V_{11}\), y su traza es \(11-1=10\). Las demás
identidades se obtienen al aplicar los proyectores a cada
sumando.
\end{proof}

La primera reducción elimina exclusivamente el modo uniforme.
La segunda elimina la dirección determinada por \(K\) en
la representación fijada. Si se cambia simultáneamente la carta
por una matriz ortogonal \(S\), el vector se transforma por
\(u_K\mapsto Su_K\) y el proyector por
\(P_{10}\mapsto SP_{10}S^{-1}\). Esta covariancia conserva
la operación geométrica, aunque cambien sus coordenadas.

\subsection{Pantalla reticular lorentziana}

Sea \(\Lambda=N_\alpha\) el retículo de Leech construido
en \ref{exc:leech}. Se considera el plano hiperbólico integral
\(U=\ZZ e_+\oplus\ZZ e_-\), con
\[
 \langle e_+,e_+\rangle=\langle e_-,e_-\rangle=0,
 \qquad\langle e_+,e_-\rangle=1.
\]
El retículo \(L_{26}=\Lambda\oplus U\) es par y
unimodular, de firma \((25,1)\). Cada vector se escribe
\(y=\lambda+ae_++be_-\). Como
\(\langle y,e_+\rangle=b\),
\begin{equation}
 e_+^\perp=\Lambda\oplus\ZZ e_+,
 \qquad e_+^\perp/\ZZ e_+\cong\Lambda.
 \label{iv:eq:26-24}
\end{equation}
El mapa del cociente es
\(q_{24}(\lambda+ae_+)=\lambda\); es sobreyectivo y tiene
núcleo \(\ZZ e_+\). La disminución de rango en dos procede
de una restricción de ortogonalidad seguida por un cociente
isotrópico. Es, por tanto, una operación de tipo diferente a
\eqref{iv:eq:p10}.

\subsection{Morfismo conjunto de pantallas}

Con el dominio \(\mathscr D_0\) de \ref{iv:def:dualidad},
sean
\[
 \mathscr G_K=\{(v_{11},P_{10}v_{11}):v_{11}\in V_{11}\},
\]
\[
 \mathscr S_{\mathrm{HMT}}=\RR^{12}\times\mathscr D_0\times e_+^\perp,
 \qquad
 \mathscr S_M^{\mathrm{alg}}=\mathscr G_K\times\mathscr D_0\times\Lambda.
\]
\begin{definition}[Mapa algebraico de pantallas coordinadas]
\label{iv:def:rm-alg}
Se define
\begin{equation}
 \mathcal R_M^{\mathrm{alg}}(v,d,y)
 =\bigl((P_{11}v,P_{10}v),d,q_{24}(y)\bigr).
 \label{iv:eq:rm-alg}
\end{equation}
\end{definition}

\begin{theorem}[Isomorfismo cocientado y entrelazamiento de dualidad]
\label{iv:thm:pantallas-isomorfas}
El mapa \eqref{iv:eq:rm-alg} induce un isomorfismo
\[
 \overline{\mathcal R}_M^{\mathrm{alg}}:
 (\RR^{12}/\RR\mathbf1)\times\mathscr D_0
 \times(e_+^\perp/\ZZ e_+)
 \longrightarrow\mathscr S_M^{\mathrm{alg}}.
\]
Entrelaza las transformaciones que actúan por \(\mathsf T_0\)
en el factor \(\mathscr D_0\) y por la identidad en los
otros dos factores. La forma \(E_{R_0}(m,w)\) se conserva
bajo esa dualidad.
\end{theorem}
\begin{proof}
La aplicación \([v]\mapsto P_{11}v\) es un isomorfismo
de \(\RR^{12}/\RR\mathbf1\) sobre \(V_{11}\).
Como \(P_{10}P_{11}=P_{10}\), el mapa
\([v]\mapsto(P_{11}v,P_{10}v)\) identifica ese cociente
con \(\mathscr G_K\); su inversa toma la primera componente
y su clase. Por \eqref{iv:eq:26-24}, \(q_{24}\) induce el
isomorfismo reticular del tercer factor. El segundo se
transporta mediante la identidad. El producto de los tres
mapas prueba el isomorfismo.

Los proyectores y \(q_{24}\) no modifican \(d\), mientras
la dualidad actúa únicamente en \(d\). Por ello las dos
composiciones son idénticas. La conservación de la forma
es el teorema \ref{iv:thm:dualidad-escalar}.
\end{proof}

El grafo \(\mathscr G_K\) conserva la coordenada anterior
\(v_{11}\) junto con su proyección. El teorema demuestra
la equivalencia de las presentaciones cocientadas; la reducción
aislada \(V_{11}\to V_{10}\) conserva su núcleo
\(\ell_K\). Esta precisión establece qué información
retiene el isomorfismo y qué información elimina cada
proyección que lo compone.

\subsection{Coordinación con la realización excepcional}

Sobre el dominio \(\mathcal X_{\mathrm{HMT}}\), con fase, orientación,
cociente, frontera, incidencia, gradación y memoria, las lecturas
\(\varepsilon_{\mathrm{exc}}\) y \(\varepsilon_{\mathrm{scr}}\) extraen
los datos excepcionales y de pantalla. El corredor construye el álgebra
graduada \(\mathfrak V(x)\cong V^\natural\), con su producto de vértices
y sus automorfismos; \(\mathfrak S_M(x)\) reúne las pantallas, cocientes
y operadores anteriores. La asignación conjunta
\(x\mapsto(\mathfrak V(x),\mathfrak S_M(x))\) tiene, por tanto, un
álgebra como primer componente y una familia de espacios y mapas como segundo.
La acción del Monstruo pertenece al primero; la reducción dirigida por \(K\)
y la dualidad compacta, al segundo. Las dos lecturas fijan la procedencia
común y cada construcción conserva su tipo.
El isomorfismo de esta sección recupera el teorema de
pantallas del integral con sus factores y su prueba
explícitos~\cite{hmtedicionintegral}.

\clearpage
\input{sections/04_accion_realizacion.tex}
\clearpage
\section{Conclusiones}
\label{iv:conclusiones}

El recorrido desarrollado enlaza la generación aritmética, la incidencia
excepcional y la realización dimensional mediante operaciones
especificadas. Su antecedente es el estado APP--TRIT--TPK con
orientación, hojas, cocientes, frontera y memoria. Las coordenadas
numéricas y los objetos excepcionales resultan de representaciones
de ese estado. Conservar el antecedente permite explicar conjuntamente
qué se evalúa y qué estructura hace posible la evaluación.

El registro \(K\) proporciona un resultado central de esta exposición.
Su construcción dodecafásica precede tanto a la lectura racional de
sus bloques como a la selección de incidencia y a la dirección
\(u_K\). En la representación real fijada, la positividad de
\(\|u_K\|^2\) determina una recta y el proyector de rango \(10\).
En el corredor excepcional, los soportes y la bandera seleccionada
conducen a la construcción del retículo y a la realización de
\(V^\natural\). La misma procedencia no hace idénticos estos
codominios: proporciona los mapas concretos que permiten compararlos
conservando sus acciones y la información que transmiten.

La representación de fase añade una consecuencia operatoria precisa
a la conservación de la unidad. La prolongación unital de todas las
fases conserva la traza normalizada y produce el cierre UHF nonádico.
Su representación tracial es un factor hiperfinito de tipo \(II_1\).
Los rangos finitos normalizados admiten en él todas las dimensiones
traciales del intervalo unidad. La continuación marcada tiene otra
inclusión y otro cierre. Este contraste identifica la operación que
conserva la normalización global y muestra por qué el dato de
refinamiento forma parte del resultado.

La dualidad de radio tiene aquí una formulación escalar y otra
operatoria, unidas por la misma involución. La igualdad de formas
cuadráticas se prolonga a una conjugación unitaria de Hamiltonianos
autoadjuntos con sus dominios. El cambio de sección residual requiere
transportar la relación \(L_9\); sobre ese cociente, la evaluación
mínima y la dualidad covariante están bien definidas. Las identidades
se mantienen para radios recíprocos y admiten la normalización de
escala de cuerda expuesta en el texto.

Las reducciones \(12\to11\to10\) y \(26\to24\) conservan,
respectivamente, la dirección determinada por \(K\) y la
presentación isotrópica del retículo. El morfismo conjunto de pantallas
coordina ambas con el sector compacto. Su prueba identifica los
cocientes, los núcleos y las transformaciones entrelazadas. La
perforación ternaria añade una realización combinatoria con su
distancia y su partición perfecta, cuyo rango queda explicado por
la operación sobre el código.

La acción sobre trayectorias se construye conservando el dato
precedente en cada composición. Los cambios de dirección y de tipo
aritmético forman entonces un cociclo exactamente aditivo. Ese
resultado fija una propiedad interna del transporte genealógico.
La superficie de mundo, la acción de Polyakov y los campos
undecadimensionales ocupan la etapa de realización: su compatibilidad
se expresa mediante los mapas y las condiciones operatorias
declarados. El teorema de transporte de la supercarga conserva
\(Q^2=H\) sobre la imagen correspondiente y precisa los dominios
necesarios para esa igualdad.

El contenido demostrado reúne, por tanto, la incidencia que alcanza
Moonshine, la dirección dodecafásica, la realización hilbertiana,
las pantallas, la dualidad compacta y la composición de acción.
La identificación física completa conserva los requisitos de sus
campos, tensiones, osciladores, amplitudes y límites declarados.
La distinción permite expresar con exactitud la contribución de cada
construcción: las pruebas algebraicas proporcionan una estructura
determinada y la realización física establece qué dinámica conserva
esa estructura. Este es el criterio de continuidad que enlaza el
presente artículo con los desarrollos posteriores de estadística,
radiación y estructura de las especies.

\clearpage
\begin{thebibliography}{99}
\phantomsection
\addcontentsline{toc}{section}{Referencias}
\bibitem{hmtintegral} O. Haidara Fall y R. Ramos Balsa, \emph{El cierre holográfico del infinito: la estructura discreta del continuo y el origen de las constantes fundamentales}. Compilación integral de trabajo de 2\,249 páginas, conservada en la revisión editorial del 5 de septiembre de 2026, y desarrollos especializados posteriores identificados en el registro de procedencia. Manuscrito inédito. La edición de 2\,084 páginas del 1 de septiembre se conserva como antecedente documental.
\bibitem{hmtsintesis} O. Haidara Fall y R. Ramos Balsa, \emph{Holografía Modular Triádica: la estructura discreta del continuo y el origen de las constantes fundamentales}. Síntesis académica autónoma, 399 páginas, conservada en la revisión editorial del 5 de septiembre de 2026. Manuscrito inédito; la edición de refuerzo TPK de 141 páginas del 1 de septiembre se conserva como antecedente.
\bibitem{hmtnuclear} O. Haidara Fall y R. Ramos Balsa, \emph{El cierre holográfico del infinito. Holografía Modular Triádica: la estructura discreta del continuo y el origen de las constantes fundamentales}. Artículo nuclear, revisión científica rectificada del 23 de julio de 2026, 97 páginas. Manuscrito inédito; antecedente histórico conservado para la procedencia de las construcciones.
\bibitem{hmtholonomia} O. Haidara Fall y R. Ramos Balsa, \emph{Doble proyección holográfica de un cristal temporal aperiódico}. Monografía de investigación, 2026, capítulos sobre aritmética de historias, holonomía y doble realización. Se consultan las versiones históricas de 65 y 775 páginas como fuentes de desarrollos especializados; la procedencia de las incorporaciones se conserva en el registro de esta revisión.
\bibitem{hmtdesarrollos} O. Haidara Fall y R. Ramos Balsa, \emph{Desarrollo matemático de la holonomía nonádica y del cristal aperiódico}; \emph{Álgebra holonómica universal y Euler trítico}. Documentos de investigación en elaboración desde el 3 de septiembre de 2026, con sus registros de procedencia y verificaciones algebraicas. Se utilizan los desarrollos individualizados de aritmética de trayectorias, vacancia central y realizaciones cuadráticas.
\bibitem{feynman1948} R. P. Feynman, ``Space-Time Approach to Non-Relativistic Quantum Mechanics'', \emph{Reviews of Modern Physics} \textbf{20}, 367--387 (1948). \href{https://doi.org/10.1103/RevModPhys.20.367}{doi:10.1103/RevModPhys.20.367}.
\bibitem{feynman1949} R. P. Feynman, ``The Theory of Positrons'', \emph{Physical Review} \textbf{76}, 749--759 (1949). \href{https://doi.org/10.1103/PhysRev.76.749}{doi:10.1103/PhysRev.76.749}.
\bibitem{feynman1965} R. P. Feynman, ``The Development of the Space-Time View of Quantum Electrodynamics'', conferencia Nobel, 11 de diciembre de 1965. \href{https://www.nobelprize.org/prizes/physics/1965/feynman/lecture/}{Texto de la conferencia}.
\bibitem{mustovicary} B. Musto y J. Vicary, ``Quantum Latin Squares and Unitary Error Bases'', \emph{Quantum Information and Computation} \textbf{16}, 1318--1332 (2016). \href{https://doi.org/10.26421/QIC16.15-16-4}{doi:10.26421/QIC16.15-16-4}.
\bibitem{delascuevas2020} G. De las Cuevas, T. Drescher y T. Netzer, ``Quantum Magic Squares: Dilations and Their Limitations'', \emph{Journal of Mathematical Physics} \textbf{61}, 111704 (2020). \href{https://doi.org/10.1063/5.0022344}{doi:10.1063/5.0022344}.
\bibitem{delascuevas2023} G. De las Cuevas, T. Netzer e I. Valentiner-Branth, ``Magic Squares: Latin, Semiclassical, and Quantum'', \emph{Journal of Mathematical Physics} \textbf{64}, 022201 (2023). \href{https://doi.org/10.1063/5.0127393}{doi:10.1063/5.0127393}.
\bibitem{dirac1933} P. A. M. Dirac, ``The Lagrangian in Quantum Mechanics'', \emph{Physikalische Zeitschrift der Sowjetunion} \textbf{3}, 64--72 (1933). \href{https://www.informationphilosopher.com/solutions/scientists/dirac/Lagrangian_1933.pdf}{Reproducción del artículo original}.
\bibitem{pauli1946} W. Pauli, ``Exclusion Principle and Quantum Mechanics'', conferencia Nobel, 13 de diciembre de 1946, en \emph{Nobel Lectures, Physics 1942--1962}, Elsevier, 1964, pp. 27--43, especialmente p. 42. \href{https://www.nobelprize.org/uploads/2018/06/pauli-lecture.pdf}{Texto de la conferencia}.
\bibitem{codata2018} E. Tiesinga, P. J. Mohr, D. B. Newell y B. N. Taylor, ``CODATA Recommended Values of the Fundamental Physical Constants: 2018'', \emph{Reviews of Modern Physics} \textbf{93}, 025010 (2021). \href{https://doi.org/10.1103/RevModPhys.93.025010}{doi:10.1103/RevModPhys.93.025010}. \href{https://physics.nist.gov/cuu/Constants/ArchiveASCII/allascii_2018.txt}{Tablas archivadas de 2018}.
\bibitem{codata2022} P. J. Mohr, D. B. Newell, B. N. Taylor y E. Tiesinga, ``CODATA Recommended Values of the Fundamental Physical Constants: 2022'', \emph{Reviews of Modern Physics} \textbf{97}, 025002 (2025). \href{https://doi.org/10.1103/RevModPhys.97.025002}{doi:10.1103/RevModPhys.97.025002}. \href{https://physics.nist.gov/cuu/Constants/Table/allascii.txt}{Tabla de constantes}.
\bibitem{flm} I. B. Frenkel, J. Lepowsky y A. Meurman, \emph{Vertex Operator Algebras and the Monster}, Pure and Applied Mathematics, vol. 134, Academic Press, 1988.
\bibitem{borcherds} R. E. Borcherds, ``Monstrous Moonshine and Monstrous Lie Superalgebras'', \emph{Inventiones Mathematicae} \textbf{109}, 405--444 (1992). \href{https://doi.org/10.1007/BF01232032}{doi:10.1007/BF01232032}.
\bibitem{conwaysloane} J. H. Conway y N. J. A. Sloane, \emph{Sphere Packings, Lattices and Groups}, 3.\textsuperscript{a} ed., Grundlehren der mathematischen Wissenschaften 290, Springer, 1999. \href{https://doi.org/10.1007/978-1-4757-6568-7}{doi:10.1007/978-1-4757-6568-7}.
\bibitem{bipm2018} Conférence générale des poids et mesures, \emph{Resolution 1 of the 26th CGPM: On the revision of the International System of Units}, 2018. \href{https://www.bipm.org/en/committees/cg/cgpm/26-2018/resolution-1}{Resolución oficial}. \href{https://doi.org/10.59161/CGPM2018RES1E}{doi:10.59161/CGPM2018RES1E}.
\bibitem{dirac1931} P. A. M. Dirac, ``Quantised singularities in the electromagnetic field'', \emph{Proceedings of the Royal Society of London. Series A} \textbf{133}, 60--72 (1931). \href{https://doi.org/10.1098/rspa.1931.0130}{doi:10.1098/rspa.1931.0130}.
\bibitem{lep2006} ALEPH, DELPHI, L3, OPAL y SLD Collaborations; LEP Electroweak Working Group; SLD Electroweak Group y SLD Heavy Flavour Group, ``Precision electroweak measurements on the Z resonance'', \emph{Physics Reports} \textbf{427}, 257--454 (2006). \href{https://doi.org/10.1016/j.physrep.2005.12.006}{doi:10.1016/j.physrep.2005.12.006}.
\bibitem{hmtedicionintegral} O. Haidara Fall y R. Ramos Balsa, \emph{El cierre holográfico del infinito}. Compilación integral de trabajo conservada en la revisión editorial del 5 de septiembre de 2026, 2\,249 páginas. Manuscrito inédito; desarrollo posterior consultado junto a la edición rectora y sus fuentes.
\bibitem{hmtedicionsintesis} O. Haidara Fall y R. Ramos Balsa, \emph{Síntesis académica autónoma de Holografía Modular Triádica y Mecánica Dimensional}. Compilación del expediente de trabajo del 3 de septiembre de 2026, 399 páginas, conservada en la revisión editorial del 5 de septiembre. Manuscrito inédito.
\bibitem{hmtedicionnuclear} O. Haidara Fall y R. Ramos Balsa, \emph{Cadena compacta y rectificación probatoria de Holografía Modular Triádica}. Compilación de trabajo del 3 de septiembre de 2026, 144 páginas, conservada en la revisión editorial del 5 de septiembre. Manuscrito inédito.
\bibitem{hmtrecopilacion} O. Haidara Fall y R. Ramos Balsa, \emph{Borrador de materiales para la generación de \(\pi,\varphi,e\), la constante de estructura fina y el electrón}. Recopilación preparatoria, 6 de septiembre de 2026, 827 páginas. Documento de trabajo y procedencia del presente artículo.
\bibitem{chenlamshimakura2016} H.-Y. Chen, C. H. Lam y H. Shimakura, ``\(\mathbb Z_3\)-orbifold construction of the Moonshine vertex operator algebra and some maximal \(3\)-local subgroups of the Monster'', prepublicación, arXiv:1606.05961 (2016), versión 2 (2017). \href{https://arxiv.org/abs/1606.05961}{arXiv:1606.05961}.
\bibitem{abelamyamada2017} T. Abe, C. H. Lam y H. Yamada, ``A remark on \(\mathbb Z_p\)-orbifold constructions of the Moonshine vertex operator algebra'', prepublicación, arXiv:1705.09022 (2017). \href{https://arxiv.org/abs/1705.09022}{arXiv:1705.09022}.
\bibitem{kmconwaynorton1979} J. H. Conway y S. P. Norton, ``Monstrous Moonshine'', \emph{Bulletin of the London Mathematical Society} \textbf{11}, 308--339 (1979). \href{https://doi.org/10.1112/blms/11.3.308}{doi:10.1112/blms/11.3.308}.
\end{thebibliography}

\end{document}
